% ============================================================================
% SECTION V — CONCLUSION AND FUTURE WORK      target 3–4 pages, 1300–1800 words
%
% Winners write LONG, STRUCTURED, QUANTITATIVE conclusions. The 2025 Chemistry
% Silver's conclusion runs ~1.5 pages and restates each result numerically rather
% than gesturing at "success". Mimic that.
%
% TRIPLE-ANCHORING: every headline number here must appear IDENTICALLY in the
% abstract and in Results. scripts/check_numbers.py enforces this.
% ============================================================================
\section{Conclusion and future work}
\label{sec:conclusion}

% ==========================================================================
% CHECKLIST — Section V
%   [ ] Restate EVERY headline number, identical to abstract and Results
%   [ ] Recall the five Findings BY NUMBER
%   [ ] One unhedged paragraph for the unified mechanistic statement
%   [ ] Limitations specific and unembarrassed, not a ritual disclaimer
%   [ ] Transferability subsection — this is the criterion C5 payload
% ==========================================================================

\subsection{Summary of findings}
\label{sec:summary}
% FIVE NUMBERED BLOCKS:
%  1. WHAT WAS MADE AND CONFIRMED — material, loading, characterisation evidence,
%     one sentence each. Note honestly that the ossein was commercially sourced.
%  2. WHAT WAS MEASURED — all headline numbers restated: capacities in BOTH units,
%     removals, alpha values against the control, thermodynamic parameters, column
%     performance, regeneration retention.
%  3. WHAT WAS COMPUTED AND WHAT IT EXPLAINED — ordering reproduced; electrostatic
%     model falsified; f_orb ordering; charge-transfer channel; hemidirection
%     descriptors; desolvation penalties.
%  4. Recall Findings 1–5 BY NUMBER.

\subsection{A unified mechanistic statement}
\label{sec:unified}
% ONE PARAGRAPH, NO HEDGING. The claim in its defensible form: selectivity for
% lead at a galloyl pocket is a synergistic coordinative effect of orbital
% covalency, hemidirected accommodation of the 6s2 lone pair, and a reduced
% desolvation penalty — NOT a charge-density effect.
%
% BE PRECISE ABOUT WHAT KIND OF SELECTIVITY: the sorbent captures a larger
% FRACTION of the lead present at a stated molar deficit. On a molar basis it does
% not take up more lead than copper. Saying this plainly in the conclusion is what
% separates a report that understands its own result from one that does not.

\subsection{Limitations}
\label{sec:limitations}
% SPECIFIC AND UNEMBARRASSED. Every item here is already known from the audit:
%  - No sorbent-free ternary blank was run, and both ternary compositions compute
%    supersaturated with respect to anglesite. Precipitation cannot be excluded by
%    measurement, only by calculation and by the contemporaneous observation that
%    the solutions remained clear.
%  - Every sample was filtered at 0.45 um, so any metal removed as a solid would be
%    counted as sorption.
%  - TGA and nitrogen physisorption were unavailable; no surface area or pore
%    structure is reported and capacity is reported per unit mass.
%  - The cluster model ignores the polymer matrix, site heterogeneity and explicit
%    second-shell water; implicit solvation cannot capture specific hydrogen
%    bonding into the vacant hemisphere — precisely the region the mechanism claims
%    is empty.
%  - The computed and measured free-energy differences agree in ordering but not in
%    magnitude, for the intrinsic-site versus ensemble reasons given in §4.6.3.
%  - Only one galloyl unit is modelled; a bis-galloyl pocket was not computed.
%  - Only three regeneration cycles; no test against real effluent, hardness ions,
%    organic matter, or trace (ug/L) lead.
%  - No binary systems were run.
%  - Ambient temperature was maintained at 25 C but the report does not claim a
%    measured mean.
%  - The relative contribution of covalent grafting versus strong physisorption is
%    not resolved by the measurements made here.

\subsection{Future work}
\label{sec:future}
%  - Repeat the competition with Cu(NO3)2 and Zn(NO3)2 so that all three metals
%    share a non-complexing counter-ion, removing the anglesite ambiguity entirely.
%  - Sorbent-free blanks at every competitive condition.
%  - Bis-galloyl pocket model; explicit second solvation shell; QM/MM or a periodic
%    model of the fibre surface.
%  - Free-energy perturbation or a site-distribution model to address the magnitude
%    discrepancy.
%  - Binary systems; hardness-ion matrix; real effluent; trace-level lead.
%  - Long-term column stability beyond three cycles; leaching over 20+ cycles.

\subsection{Transferability of the design principle}
\label{sec:transferability}
% THE CRITERION C5 PAYLOAD. Argue that the orbital-covalency and desolvation-
% penalty criterion is a GENERAL DESIGN DESCRIPTOR for separating soft,
% polarisable, lone-pair-active ions — Pb(II), Tl(I), Bi(III), Sb(III), Sn(II) —
% from harder transition ions at oxygen-donor sites.
% Plus the waste-valorisation angle and the demonstrated concentration factor for
% actual recovery rather than mere removal.
